\documentclass[10pt,a4paper]{amsart}
\usepackage[margin=19mm]{geometry}
\usepackage{amsmath,amssymb,mathtools}
\usepackage[scr=rsfs,cal=euler]{mathalfa}
\usepackage{xfrac}
\usepackage[unicode,hidelinks,bookmarks=false]{hyperref}
\newcommand{\ie}{i.e.\ }
\newcommand{\cf}{cf.\ }
\newcommand{\resp}{respectively\ }
\newcommand{\Subsection}{\subsection}
\providecommand{\nobreakdash}{\nobreak\hbox{-}}
\setlength{\emergencystretch}{2em}
\multlinegap0pt
\usepackage{mathtools}
\usepackage{amssymb}
\usepackage{xcolor}
\usepackage{algorithm}\usepackage{algpseudocode}
\usepackage{enumerate}
\usepackage{float}
\newtheorem{lemma}{Lemma}
\title{Construction of cyclic codes with large minimum distance from power functions over odd characteristic finite fields}
\author{Mrinal Kanti Bose, Abhay Kumar Singh}
\date{Source: arXiv:2606.03638v1 (CC BY 4.0)}
\begin{document}
\maketitle

\noindent\emph{Source and licence.} Construction of cyclic codes with large minimum distance from power functions over odd characteristic finite fields. Version arXiv:2606.03638v1 (CC BY 4.0), distributed by arXiv under the Creative Commons Attribution 4.0 International licence, http://creativecommons.org/licenses/by/4.0/, as recorded in the arXiv licence metadata for this exact version and shown on the abstract page https://arxiv.org/abs/2606.03638v1. Full licence notice: \texttt{licenses/arxiv-2606.03638-CC-BY-4.0.txt}.\par\smallskip

\noindent\textit{Editorial scope.} Counted unit: the article's lemma L1 (source0.tex lines 699--719). The article's complete original proof of it is delivered with it (source0.tex lines 729--755, one \texttt{proof} environment of 27 lines, which the article places away from the statement). No mathematics is written, repaired or re-derived: the statement and the proof are the article's own lines, in the article's own notation, and no retained line is substituted or re-flowed.  Retained uncounted as the source-local context the proof needs: lines 676--678.  Nothing was omitted from any retained span, so the package carries no omission record.  No retained line contains a citation, so the wrapper carries no bibliography and no bibliography entry.  No retained line contains a figure environment, an image-inclusion command or drawing code, so no image is delivered and the package declares no asset dependency.  Editorial formatting only, in addition: an AMS \texttt{amsart} preamble that restates the article's own declarations and macros used by the retained lines, namely the environments \texttt{lemma} on separate counters, together with the standard packages those lines need.  The retained environments print in this document as Lemma 1 in order of appearance, whereas the article prints the same items with the numbering of its own full document.  The complete native source of the article as distributed in the arXiv e-print for this exact version is bundled unchanged as \texttt{sources/201990/source0.tex}.
\begin{equation}\label{Eq1}
    D_{t}(q)=\{\sum_{i=0}^{t-1}a_{i}q^{i}:a_{i}\text{'s, not all zero, is either $0$ or $1$ for all $i=0,1,2,\cdots,t-1$}\}.
\end{equation}

\begin{lemma}\label{L1}
    Let $t\in\mathbb{N}$ and $j\in D_{t+1}(q)$ be an integer not divisible by $q$, where $D_{t}(q)$ is defined as in $(\ref{Eq1})$. Then 
    \begin{enumerate}
        \item[$1.$] For $j\in D_{t}(q)$, $B_{j}^{(t+1)}=B_{j}^{(t)}\cup\{jq^{\epsilon_{j}^{(t)}}\}$ and $\epsilon_{j}^{(t+1)}=\epsilon_{j}^{(t)}+1$. 
        \item[$2.$] For $j\in D_{t+1}(q)\backslash D_{t}(q)$, $B_{j}^{(t)}=\{j\}$ and $\epsilon_{j}^{(t+1)}=1$. 
    \end{enumerate}
    \begin{proof}
      From the definition of $\epsilon_{j}^{(t+1)}$, we know that 
      \begin{align*}
          \epsilon_{j}^{(t+1)} &=\begin{cases}
              1, \text{ if }j=\frac{q^{t+1}-1}{q-1} \\
              \Big\lceil\operatorname{log}_{q}\left(\frac{q^{t+1}-1}{j(q-1)}\right)\Big\rceil,\text{ if }1\leq j<\frac{q^{t+1}-1}{q-1}
          \end{cases} 
      \end{align*}
      If $j\in D_{t}(q)$ and $q\nmid j$, then $\frac{q^{t+1}-1}{(q-1)j}=\frac{q(q^t-1)}{(q-1)j}+\frac{1}{j}$ gives $\frac{q(q^t-1)}{(q-1)j}<\frac{q^{t+1}-1}{(q-1)j}\leq \frac{q(q^t-1)}{(q-1)j}+1$, implying that $1+\operatorname{log}_{q}\left(\frac{q^t-1}{(q-1)j}\right)<\operatorname{log}_{q}\left(\frac{q^{t+1}-1}{(q-1)j}\right)<2+\operatorname{log}_{q}\left(\frac{q^t-1}{(q-1)j}\right)$. Since $\operatorname{log}_{q}\left(\frac{q^t-1}{(q-1)j}\right)$ can not be a nonzero integer, we conclude that $\epsilon_{j}^{(t+1)}=\epsilon_{j}^{(t)}+1$ for all $j\in D_{t}(q)$.
      \newline
      If $j\in D_{t+1}(q)\backslash D_{t}(q)$, $q\nmid j$, note that $1+q^t\leq j\leq\frac{q^{t+1}-1}{q-1}$. For $j=\frac{q^{t+1}-1}{q-1}$, we know that $\epsilon_{j}^{(t+1)}=1$ hold by the definition. For $j\neq \frac{q^{t+1}-1}{q-1}$, $j<\frac{q^{t+1}-1}{q-1}<qj$ would give $0<\operatorname{log}_{q}\left(\frac{q^{t+1}-1}{j(q-1)}\right)<1$. Therefore, $\epsilon_{j}^{(t+1)}=1$ for all $j\in D_{t+1}(q)\backslash D_{t}(q)$.
      \vskip 1pt
      Thus, the proof of the lemma follows from the the definition of $B_{j}^{(t)}$.
    \end{proof}
\end{lemma}

    \begin{proof}
        For $t=1$, we have $D_{t}(q)=\{1\}$. Observe that $\epsilon_{1}^{(1)}=1$ due to the definition. For all $t\geq 2$, since $q^{t-1}<\frac{q^t-1}{q-1}<q^t$, we obtain $\epsilon_{1}^{(t)}=\Big\lceil\operatorname{log}_{q}\left(\frac{q^t-1}{q-1}\right)\Big\rceil=t$.
        %\newline
       \par Take $t=2$ and $j\in D_{2}(q)$ with $q\nmid j$. Note that $D_{1}(q)=\{1\}$ and $D_{2}(q)=\{1,q,1+q\}$. Then, from Lemma \ref{L1}, we obtain
       \begin{align*}
           \epsilon_{j}^{(t)}=\epsilon_{j}^{(2)} &=\begin{cases}
               \epsilon_{j}^{(1)}+1=2;\text{ for }j=1, \\
               1;\text{ for }j\in D_{2}(q)\backslash D_{1}(q)\text{ and }q\nmid j.
           \end{cases}
       \end{align*}
      % $\epsilon_{1}^{(2)}=\epsilon_{1}^{(1)}+1=2$ and $\epsilon_{1+q}^{(2)}=1$. 
       \par Take $t=3$ and $j\in D_{3}(q)$ with $q\nmid j$. Note that $D_{2}(q)=\{1,q,1+q\}$ and $D_{3}(q)=\{1,q,q^2,1+q,1+q^2,q+q^2,1+q+q^2\}$. Then from Lemma \ref{L1}, we obtain
        \begin{align*}
           \epsilon_{j}^{(t)}=\epsilon_{j}^{(3)} &=\begin{cases}
               \epsilon_{j}^{(2)}+1;\text{ for }j\in D_{2}(q)\text{ and }q\nmid j, \\
               1;\text{ for }j\in D_{3}(q)\backslash D_{2}(q)\text{ and }q\nmid j
           \end{cases}
           =\begin{cases}
               3;\text{ for }j=1, \\
               2;\text{ for }j\in D_{2}(q)\backslash D_{1}(q)\text{ and }q\nmid j, \\
               1;\text{ for }j\in D_{3}(q)\backslash D_{2}(q)\text{ and }q\nmid j.
           \end{cases}
       \end{align*}
      % $\epsilon_{1}^{(3)}=\epsilon_{1}^{(2)}+1=3$, $\epsilon_{1+q}^{(3)}=\epsilon_{1+q}^{(2)}+1=2$ and $\epsilon_{1+q^2}^{(3)}=\epsilon_{1+q+q^2}^{(3)}=1$.
       \vskip 1pt
     \par  By continuing this reasoning for all values of $t$, we obtain the desired result.
    \end{proof}

\end{document}
